\chapter{Bộ truy vấn đánh giá}
\label{app:truy-van}

Bảng~\ref{tab:pl-truy-van} liệt kê nguyên văn 155 truy vấn dùng ở Chương~\ref{chap:danh_gia}, giữ cách viết của người soạn, kể cả từ mượn tiếng Anh. Nhóm được định nghĩa ở Mục~\ref{sec:du-lieu-cach-do}: A là một thuộc tính, B là nhiều thuộc tính, C là từ địa phương hoặc từ mượn tiếng Anh, D là theo phong cách. Truy vấn số 1 đến 133 là lô thứ nhất, xếp theo nhóm; số 134 đến 155 là lô thứ hai, cũng chia theo bốn nhóm đó.

\begingroup
\scriptsize
\setlength{\tabcolsep}{3pt}
\renewcommand{\arraystretch}{0.9}
\begin{longtable}{@{}r c p{5.7cm} r c p{5.7cm}@{}}
\caption{Bộ 155 truy vấn đánh giá}\label{tab:pl-truy-van}\\
\toprule
\textbf{Mã} & \textbf{Nhóm} & \textbf{Truy vấn} & \textbf{Mã} & \textbf{Nhóm} & \textbf{Truy vấn} \\
\midrule
\endfirsthead
\toprule
\textbf{Mã} & \textbf{Nhóm} & \textbf{Truy vấn} & \textbf{Mã} & \textbf{Nhóm} & \textbf{Truy vấn} \\
\midrule
\endhead
\bottomrule
\endfoot
1 & A & áo sơ mi nữ trắng & 79 & C & đầm y2k dáng ôm màu hồng \\
2 & A & áo thun cổ tròn đen & 80 & C & áo thun basic form rộng unisex \\
3 & A & áo len dệt kim xám & 81 & C & quần culottes ống loe kiểu Hàn \\
4 & A & áo cardigan màu kem & 82 & C & áo cardigan mỏng nhẹ mùa hè \\
5 & A & áo khoác jean bạc màu & 83 & C & váy tomboy cá tính năng động \\
6 & A & áo gile len không tay & 84 & C & áo polo nữ dáng rộng basic \\
7 & A & quần tây đen ống suông & 85 & C & quần jean rách gối cá tính \\
8 & A & quần short denim xanh & 86 & C & áo dệt kim knit dáng crop ngắn \\
9 & A & chân váy xếp ly đen & 87 & C & set đồ tây công sở nữ basic \\
10 & A & áo khoác dạ màu nâu & 88 & C & áo turtleneck cổ lọ dệt kim \\
11 & A & đầm suông màu đỏ & 89 & C & quần wide leg ống suông rộng \\
12 & A & áo liền quần màu trắng & 90 & C & áo croptop rib bo sát người \\
13 & A & áo choàng không tay đen & 91 & C & váy midi dáng A-line \\
14 & A & kính mắt gọng tròn & 92 & C & áo blazer oversize dáng vintage \\
15 & A & mũ lưỡi trai đen & 93 & C & quần jean skinny bó sát \\
16 & A & băng đô vải nhung & 94 & C & áo sweater knit chunky dày \\
17 & A & cà vạt kẻ sọc & 95 & C & đầm slip dress dáng suông \\
18 & A & găng tay da đen & 96 & C & áo bomber jacket dáng ngắn \\
19 & A & đồng hồ dây da nâu & 97 & C & quần cargo nhiều túi hộp \\
20 & A & dây nịt da đen & 98 & C & áo tank top ba lỗ basic \\
21 & A & quần bó tối màu & 99 & C & giày sneaker trắng basic \\
22 & A & vớ cổ cao trắng & 100 & C & túi tote vải canvas basic \\
23 & A & giày cao gót đen & 101 & C & mũ bucket hat vải kaki \\
24 & A & túi xách da nâu & 102 & C & áo vest oversize kiểu boyfriend \\
25 & A & khăn choàng len dài & 103 & C & chân váy pleated skirt xếp ly \\
26 & A & dù gấp gọn nhỏ & 104 & D & đồ đi làm phong cách công sở tối giản \\
27 & A & áo phông tay ngắn nam & 105 & D & outfit phong cách Hàn Quốc trẻ trung \\
28 & A & quần jean rách nhẹ & 106 & D & đồ dạo phố năng động mùa hè \\
29 & A & đầm hai dây màu đen & 107 & D & trang phục vintage thập niên 90 \\
30 & A & áo khoác gió mỏng & 108 & D & đồ đi biển thoải mái mùa hè \\
31 & A & chân váy chữ A ngắn & 109 & D & set đồ dự tiệc sang trọng buổi tối \\
32 & A & áo sơ mi kẻ caro nam & 110 & D & đồ mặc nhà thoải mái rộng rãi \\
33 & A & quần jogger thể thao & 111 & D & phong cách đường phố cá tính unisex \\
34 & A & áo len cổ lọ dày & 112 & D & trang phục thanh lịch đi làm mùa đông \\
35 & A & ví cầm tay nhỏ & 113 & D & đồ đi học năng động trẻ trung \\
36 & A & giày thể thao trắng & 114 & D & phong cách tối giản màu trung tính \\
37 & A & áo vest nữ đen & 115 & D & outfit dạo phố phong cách Nhật \\
38 & A & quần tây kẻ sọc & 116 & D & đồ mặc đi tiệc sinh nhật nữ tính \\
39 & A & áo khoác lông giả & 117 & D & trang phục đi phỏng vấn lịch sự \\
40 & A & khăn quàng cổ lụa & 118 & D & đồ mặc đi làm mùa thu ấm áp \\
41 & B & áo sơ mi lụa tay dài công sở màu trắng ngà & 119 & D & outfit đi chơi cuối tuần năng động \\
42 & B & áo khoác bomber nam màu xanh rêu vải dù & 120 & D & phong cách nữ tính nhẹ nhàng \\
43 & B & quần jean ống rộng lưng cao màu xanh nhạt & 121 & D & trang phục dự đám cưới trang trọng \\
44 & B & đầm hai dây satin màu hồng đi dự tiệc & 122 & D & đồ mặc đi du lịch tiện dụng \\
45 & B & áo len cổ tim mỏng màu than mặc mùa thu & 123 & D & phong cách cá tính nổi bật \\
46 & B & quần tây ống đứng vải tuyết mưa màu đen công sở & 124 & D & outfit công sở mùa hè thoáng mát \\
47 & B & áo khoác dạ dáng dài màu camel cổ bẻ & 125 & D & trang phục đi làm phong cách hiện đại \\
48 & B & chân váy chữ A hoa nhỏ tay bồng đi dự tiệc & 126 & D & đồ mặc đi chơi lễ hội âm nhạc \\
49 & B & áo phông cotton oversize màu trắng unisex & 127 & D & phong cách thanh lịch quý phái \\
50 & B & chân váy da bóng dáng bút chì màu đen dài gối & 128 & D & outfit đi cà phê cuối tuần đơn giản \\
51 & B & áo len dệt kim màu hồng pastel cổ tròn ngắn tay & 129 & D & trang phục thể thao năng động hàng ngày \\
52 & B & quần jogger thể thao nữ màu xám lưng rút & 130 & D & đồ mặc mùa đông ấm áp mà vẫn thời trang \\
53 & B & áo blazer nữ dáng ngắn vải tweed màu đen & 131 & D & phong cách tomboy cá tính mạnh mẽ \\
54 & B & đầm suông linen màu nude đi làm mùa hè & 132 & D & outfit đi làm freelancer thoải mái \\
55 & B & áo khoác denim nữ dáng oversize màu xanh nhạt & 133 & D & trang phục dạo phố mùa xuân nhẹ nhàng \\
56 & B & quần short kaki túi hộp màu xanh rêu & 134 & A & áo thun cổ tim nữ trắng \\
57 & B & đầm dạ hội đen dài tay ren dự tiệc tối & 135 & A & quần âu nam màu xanh navy \\
58 & B & áo sơ mi nam trắng tay ngắn công sở vải cotton & 136 & A & áo khoác kaki màu vàng đất \\
59 & B & chân váy midi xếp ly màu beige đi làm & 137 & A & chân váy da lộn màu nâu \\
60 & B & áo cardigan mỏng màu pastel cài nút gỗ & 138 & A & áo len tay dài màu tím \\
61 & B & quần culottes ống loe màu đen vải kaki & 139 & A & đầm ren màu trắng ngà \\
62 & B & áo khoác phao dày màu đen đi mùa đông & 140 & B & áo khoác nỉ có mũ màu đen \\
63 & B & áo len cổ lọ dệt kim màu rêu ôm dáng & 141 & B & quần short jean lưng cao ống loe nhẹ \\
64 & B & váy hai dây lụa satin màu vàng ánh kim & 142 & B & áo sơ mi voan hoa nhỏ tay phồng \\
65 & B & áo khoác vest nữ dáng dài vải nỉ màu xám & 143 & B & đầm babydoll tay bồng màu xanh mint \\
66 & B & quần tây baggy màu đen công sở nam & 144 & B & áo khoác gió chống nước màu cam \\
67 & B & áo croptop len dệt kim màu trắng ngắn tay & 145 & B & quần legging tập gym màu đen bó sát \\
68 & B & áo thun cotton dày dặn màu đen form rộng & 146 & C & áo vest nam form rộng dáng Hàn \\
69 & B & đầm body ôm sát màu đen dự tiệc tối & 147 & C & quần jean mom fit lưng cao \\
70 & B & áo khoác cardigan dài vải len màu nâu trà & 148 & C & đầm maxi dáng dài thướt tha \\
71 & B & giày cao gót mũi nhọn màu đen da bóng & 149 & C & áo peplum eo nhấn dáng xoè \\
72 & B & túi xách nhỏ đeo chéo màu đen da mềm & 150 & C & quần palazzo ống rộng bay bổng \\
73 & B & mũ len bo tròn màu xám đội mùa đông & 151 & D & trang phục đi làm remote thoải mái tại nhà \\
74 & B & kính mắt gọng vuông màu đen dáng cổ điển & 152 & D & outfit đi xem phim cuối tuần đơn giản \\
75 & B & khăn choàng len dày màu đỏ đô mùa lạnh & 153 & D & phong cách preppy trẻ trung học đường \\
76 & C & áo croptop ulzzang màu trắng & 154 & D & trang phục chụp ảnh kỷ yếu nữ tính \\
77 & C & quần baggy ống rộng màu đen & 155 & D & đồ mặc leo núi dã ngoại tiện dụng \\
78 & C & áo hoodie oversize màu xám &  &  &  \\
\end{longtable}
\endgroup


\chapter{Đối soát nhãn chấm 0 tự động}
\label{app:doi-soat}

Bảng~\ref{tab:pl-doi-soat} liệt kê 35 truy vấn bị đánh dấu khi đối soát nhãn 0 tự động (định nghĩa bộ lọc ở Mục~\ref{sec:du-lieu-cach-do}). Từ khoá bộ lọc là tên danh mục mà bộ lọc chọn từ câu truy vấn. Cột Loại chia nguyên nhân bị đánh dấu thành ba nhóm: bộ lọc chọn nhầm danh mục (từ khoá giả), truy vấn phụ kiện, và truy vấn có chữ ``đầm''. Ví dụ nhóm đầu: ``Quần dài'' được chọn cho chữ ``tay dài''. Với truy vấn phụ kiện (kính, túi, giày), ảnh có đúng món được hỏi vẫn bị gán 0 vì danh mục chính của ảnh là món khác. Với truy vấn có chữ ``đầm'', bộ lọc chọn đúng loại sản phẩm; truy vấn bị đánh dấu vì câu không chứa chữ ``váy'' trong tên danh mục ``Váy đầm''.

Cột Cặp bị chấm 0 đếm theo hai cách. Với truy vấn từ khoá giả và từ khoá đầm, cột này tính mọi cặp bị bộ lọc gán 0. Với truy vấn phụ kiện, cột chỉ tính các cặp có món được hỏi trong tiêu đề ảnh. Cột Top-10 tự động đếm, trong 10 kết quả đầu của hệ DoRA, số kết quả bị bộ lọc gán nhãn 0 và chưa có người xem lại. Ở 14 truy vấn phụ kiện, cột này là 9 hoặc 10, nên Recall@10 của chúng bằng 0 do cách chấm. Ở 9 truy vấn đầm, cột này chỉ từ 0 đến 2 và hệ DoRA không trượt truy vấn nào. Trong 12 truy vấn từ khoá giả, 4 truy vấn (số 41, 57, 64 và 148) là những truy vấn hệ DoRA không tìm đúng.

\begingroup
\footnotesize
\setlength{\tabcolsep}{4pt}
\renewcommand{\arraystretch}{0.95}
\begin{longtable}{@{}r >{\raggedright\arraybackslash}p{5.2cm} >{\raggedright\arraybackslash}p{2.3cm} >{\raggedright\arraybackslash}p{3.0cm} >{\raggedleft\arraybackslash}p{1.4cm} >{\raggedleft\arraybackslash}p{1.5cm}@{}}
\caption{35 truy vấn bị đánh dấu khi đối soát nhãn 0 tự động}\label{tab:pl-doi-soat}\\
\toprule
\textbf{Mã} & \textbf{Truy vấn} & \textbf{Từ khoá bộ lọc} & \textbf{Loại} & \textbf{Cặp bị chấm 0} & \textbf{Top-10 tự động} \\
\midrule
\endfirsthead
\toprule
\textbf{Mã} & \textbf{Truy vấn} & \textbf{Từ khoá bộ lọc} & \textbf{Loại} & \textbf{Cặp bị chấm 0} & \textbf{Top-10 tự động} \\
\midrule
\endhead
\bottomrule
\endfoot
11 & đầm suông màu đỏ & Váy đầm & Từ khoá đầm & 16 & 1/10 \\
14 & kính mắt gọng tròn & Kính mắt & Phụ kiện & 28 & 10/10 \\
16 & băng đô vải nhung & Băng đô & Phụ kiện & 25 & 10/10 \\
17 & cà vạt kẻ sọc & Cà vạt & Phụ kiện & 21 & 10/10 \\
18 & găng tay da đen & Áo choàng không tay & Từ khoá giả & 28 & 8/10 \\
19 & đồng hồ dây da nâu & Đồng hồ & Phụ kiện & 24 & 10/10 \\
20 & dây nịt da đen & Dây nịt & Phụ kiện & 25 & 10/10 \\
23 & giày cao gót đen & Giày & Phụ kiện & 41 & 10/10 \\
24 & túi xách da nâu & Túi & Phụ kiện & 23 & 10/10 \\
25 & khăn choàng len dài & Áo len & Từ khoá giả & 20 & 4/10 \\
29 & đầm hai dây màu đen & Váy đầm & Từ khoá đầm & 20 & 1/10 \\
35 & ví cầm tay nhỏ & Áo choàng không tay & Từ khoá giả & 33 & 9/10 \\
36 & giày thể thao trắng & Giày & Phụ kiện & 35 & 10/10 \\
40 & khăn quàng cổ lụa & Khăn choàng & Phụ kiện & 22 & 9/10 \\
41 & áo sơ mi lụa tay dài công sở màu trắng ngà & Quần dài & Từ khoá giả & 40 & 10/10 \\
44 & đầm hai dây satin màu hồng đi dự tiệc & Váy đầm & Từ khoá đầm & 9 & 0/10 \\
50 & chân váy da bóng dáng bút chì màu đen dài gối & Quần dài & Từ khoá giả & 43 & 9/10 \\
54 & đầm suông linen màu nude đi làm mùa hè & Váy đầm & Từ khoá đầm & 15 & 0/10 \\
57 & đầm dạ hội đen dài tay ren dự tiệc tối & Quần dài & Từ khoá giả & 41 & 10/10 \\
64 & váy hai dây lụa satin màu vàng ánh kim & Chân váy & Từ khoá giả & 45 & 10/10 \\
69 & đầm body ôm sát màu đen dự tiệc tối & Váy đầm & Từ khoá đầm & 10 & 0/10 \\
71 & giày cao gót mũi nhọn màu đen da bóng & Giày & Phụ kiện & 39 & 10/10 \\
72 & túi xách nhỏ đeo chéo màu đen da mềm & Túi & Phụ kiện & 22 & 10/10 \\
73 & mũ len bo tròn màu xám đội mùa đông & Áo len & Từ khoá giả & 13 & 2/10 \\
74 & kính mắt gọng vuông màu đen dáng cổ điển & Kính mắt & Phụ kiện & 29 & 10/10 \\
75 & khăn choàng len dày màu đỏ đô mùa lạnh & Áo len & Từ khoá giả & 22 & 2/10 \\
79 & đầm y2k dáng ôm màu hồng & Váy đầm & Từ khoá đầm & 19 & 1/10 \\
83 & váy tomboy cá tính năng động & Chân váy & Từ khoá giả & 36 & 7/10 \\
91 & váy midi dáng A-line & Chân váy & Từ khoá giả & 37 & 7/10 \\
95 & đầm slip dress dáng suông & Váy đầm & Từ khoá đầm & 13 & 1/10 \\
99 & giày sneaker trắng basic & Giày & Phụ kiện & 39 & 10/10 \\
100 & túi tote vải canvas basic & Túi & Phụ kiện & 33 & 10/10 \\
139 & đầm ren màu trắng ngà & Váy đầm & Từ khoá đầm & 15 & 0/10 \\
143 & đầm babydoll tay bồng màu xanh mint & Váy đầm & Từ khoá đầm & 14 & 2/10 \\
148 & đầm maxi dáng dài thướt tha & Quần dài & Từ khoá giả & 35 & 9/10 \\
\end{longtable}
\endgroup

